\section*{Capítulo \ref{ch:g5:air-pressure-weather} --- Pressão do ar e o tempo}

\begin{solution}{exo:g5:air-pressure-weather:1}
A mesma bola é pesada cheia e depois murcha: só o \omterm{def:g2:air-around-us:air}{ar} saiu dela, e a
balança marca menos --- a diferença é o \omterm{def:g4:weight-and-mass:weight}{peso} do \omterm{def:g2:air-around-us:air}{ar} que escapou. O \omterm{def:g2:air-around-us:air}{ar}
é material com \omterm{def:g3:measuring-mass:mass}{massa}.
\end{solution}

\begin{solution}{exo:g5:air-pressure-weather:2}
É o empurrão do \omterm{def:g2:air-around-us:air}{ar} apertado do fundo da \omterm{def:g5:air-pressure-weather:pressure}{atmosfera} sobre toda
superfície em que ele encosta. Ela vem do \omterm{def:g4:weight-and-mass:weight}{peso} do oceano profundo de
\omterm{def:g2:air-around-us:air}{ar} acima de \omterm{def:g5:series-parallel-circuits:parallel}{nós}, e empurra de todas as direções --- para baixo, de
lado e para cima.
\end{solution}

\begin{solution}{exo:g5:air-pressure-weather:3}
Porque os empurrões chegam equilibrados: o \omterm{def:g2:air-around-us:air}{ar} do outro lado da sua
palma --- e o de dentro do seu corpo --- empurra de volta com
exatamente a mesma \omterm{def:g2:pushes-and-pulls:force}{força}. Só um empurrão \emph{desequilibrado}, como
no cartão ou na garrafa que amassa, mostra a \omterm{def:g2:pushes-and-pulls:force}{força} que tem.
\end{solution}

\begin{solution}{exo:g5:air-pressure-weather:4}
O \omterm{def:g4:weight-and-mass:weight}{peso} da água empurra o cartão para baixo; a \omterm{def:g5:air-pressure-weather:pressure}{atmosfera} empurra o
cartão para cima. A \omterm{def:g5:air-pressure-weather:pressure}{atmosfera} ganha --- o cartão fica e a água
continua presa.
\end{solution}

\begin{solution}{exo:g5:air-pressure-weather:5}
Você tira o \emph{\omterm{def:g2:air-around-us:air}{ar}} de dentro do canudinho. A \omterm{def:g5:air-pressure-weather:pressure}{atmosfera}, apertando
a superfície do suco no copo, então empurra o suco para cima pelo
canudinho livre: quem levanta é o céu.
\end{solution}

\begin{solution}{exo:g5:air-pressure-weather:6}
Apertar a ventosa contra a parede espremeu para fora o \omterm{def:g2:air-around-us:air}{ar} que estava
atrás dela. Sem nada empurrando de volta por baixo, o empurrão da
\omterm{def:g5:air-pressure-weather:pressure}{atmosfera} prende a ventosa nos azulejos.
\end{solution}

\begin{solution}{exo:g5:air-pressure-weather:7}
Ela enfraquece --- sobra menos oceano de \omterm{def:g2:air-around-us:air}{ar} acima para apertar.
Sinais: os ouvidos estalando numa subida rápida, a falta de \omterm{def:g2:air-around-us:air}{ar} no \omterm{def:g2:air-around-us:air}{ar}
rarefeito (e o chá que ferve frio demais).
\end{solution}

\begin{solution}{exo:g5:air-pressure-weather:8}
O empurrão da \omterm{def:g5:air-pressure-weather:pressure}{atmosfera} --- a \omterm{def:g5:air-pressure-weather:pressure}{pressão atmosférica}. Uma queda rápida
avisa que vem tempestade: leve o piquenique para dentro de casa.
\end{solution}

\begin{solution}{exo:g5:air-pressure-weather:9}
Lacrada no pico, a garrafa guardou por dentro o \omterm{def:g2:air-around-us:air}{ar} rarefeito da
montanha; no vale, o empurrão de fora ficou mais forte que o
empurrão preso lá dentro, e a garrafa dobrou. No caminho contrário,
uma garrafa lacrada no vale guarda \omterm{def:g2:air-around-us:air}{ar} \emph{forte} por dentro
enquanto o empurrão de fora enfraquece com a altura: a garrafa incha
esticada --- e um pacote de salgadinho lacrado faz o mesmo em toda
subida de montanha.
\end{solution}

\begin{solution}{exo:g5:air-pressure-weather:10}
O \omterm{def:g2:air-around-us:air}{ar} despeja sem parar das regiões pesadas de alta pressão na direção
das regiões leves de baixa pressão, e esse despejo é o \omterm{def:g2:air-around-us:wind}{vento}. Assim a
velha definição se completa: \omterm{def:g2:air-around-us:wind}{vento} é \omterm{def:g2:air-around-us:air}{ar} em movimento --- movido por
diferenças de \omterm{def:g5:air-pressure-weather:pressure}{pressão atmosférica}.
\end{solution}

\begin{solution}{exo:g5:air-pressure-weather:11}
Com a altura, o empurrão da \omterm{def:g5:air-pressure-weather:pressure}{atmosfera} sobre a água enfraquece, então
a fervura vem mais fácil: o \omterm{def:g4:changes-of-state:boiling-point}{ponto de ebulição} escorrega para bem
abaixo de \qty{100}{\celsius}. E, pela lei do patamar, a água
fervendo nunca consegue subir além do \omterm{def:g4:changes-of-state:boiling-point}{ponto de ebulição} dela ---
chama a mais só produz vapor. Ou seja: a água no alto do pico ferve
fria \emph{e fica presa nisso}: chá fraco, macarrão duro, física
satisfeita.
\end{solution}
